\documentclass[11pt, a4paper]{article}

% --- PAQUETES PRINCIPALES ---
\usepackage[utf8]{inputenc}
\usepackage[spanish, es-tabla]{babel}
\usepackage[margin=2.5cm]{geometry}
\usepackage{amsmath, amssymb, mathtools}
\usepackage{siunitx}
\usepackage{graphicx}
\usepackage{booktabs}
\usepackage{tabularx}
\usepackage[most]{tcolorbox}
\usepackage{listings}
\usepackage{xcolor}

% --- PAQUETES PARA GRÁFICAS Y CIRCUITOS ---
\usepackage{pgfplots}
\usepackage[american]{circuitikz}
\pgfplotsset{compat=1.18}

% --- CONFIGURACIÓN DE CÓDIGO (PYTHON) ---
\definecolor{codegreen}{rgb}{0,0.6,0}
\definecolor{codegray}{rgb}{0.5,0.5,0.5}
\definecolor{codepurple}{rgb}{0.58,0,0.82}
\definecolor{backcolour}{rgb}{0.96,0.96,0.96}
\lstset{
    backgroundcolor=\color{backcolour},   
    commentstyle=\color{codegreen},
    keywordstyle=\color{magenta},
    numberstyle=\tiny\color{codegray},
    stringstyle=\color{codepurple},
    basicstyle=\ttfamily\footnotesize,
    breaklines=true,                 
    numbers=left,                    
    numbersep=5pt,                  
    language=Python,
    literate={á}{{\'a}}1 {é}{{\'e}}1 {í}{{\'i}}1 {ó}{{\'o}}1 {ú}{{\'u}}1
             {Á}{{\'A}}1 {É}{{\'E}}1 {Í}{{\'I}}1 {Ó}{{\'O}}1 {Ú}{{\'U}}1
             {ñ}{{\~n}}1 {Ñ}{{\~N}}1
}

% --- CONFIGURACIÓN DE SIUNITX ---
\sisetup{output-decimal-marker = {,}, per-mode = symbol, inter-unit-product = \ensuremath{{}\cdot{}}}

% --- CABECERAS ---
\usepackage{fancyhdr}
\pagestyle{fancy}
\fancyhf{}
\lhead{\textbf{Circuitos Electrónicos III (IMT-221)}}
\rhead{Sesión 6: Recortadores y Sujetadores}
\cfoot{\thepage}
\renewcommand{\headrulewidth}{0.4pt}

% --- ENTORNOS PERSONALIZADOS ---
\newtcolorbox{aviso}[1][]{colback=red!5!white, colframe=red!75!black, fonttitle=\bfseries, title=#1}

\begin{document}

\begin{center}
    {\LARGE \textbf{Apuntes de Cátedra: Sesión 6}}\\[0.3cm]
    {\large \textbf{Limitadores (Clippers) y Sujetadores de Nivel (Clampers)}}\\
    \textit{Docente: M.Sc. Ing. Bernardo Quiroga Turdera}
    \vspace{0.3cm}
    \hrule
\end{center}

\vspace{0.3cm}

% ==========================================
% SECCIÓN 1: INTRODUCCIÓN Y CONTEXTO
% ==========================================
\section{Acondicionamiento de Señales Bipolares (El Reto del PFI)}
Un sensor piezoeléctrico o un transductor de corriente entrega una señal senoidal de $\pm \qty{10}{\volt}$. Si conectamos esta señal directamente al ADC de un microcontrolador STM32 (rango $\qty{0}{\volt}$ a $\qty{+3.3}{\volt}$) o ESP32 ($\qty{0}{\volt}$ a $\qty{+5.0}{\volt}$), los semiciclos negativos \textbf{destruirán la compuerta del conversor} y las sobretensiones quemarán el chip. 

Solucionamos esto con dos topologías:
\begin{enumerate}
    \item \textbf{Clampers (Sujetadores):} Desplazan la señal completa hacia la zona positiva sin alterar su forma.
    \item \textbf{Clippers (Recortadores):} Recortan los picos que exceden el límite de seguridad establecido.
\end{enumerate}

% ==========================================
% SECCIÓN 2: RECORTADORES (CLIPPERS)
% ==========================================
\section{Circuitos Recortadores (Clippers)}
Son redes que eliminan una porción de la onda por encima o por debajo de una referencia.

\begin{center}
\renewcommand{\arraystretch}{1.5}
\begin{tabularx}{\textwidth}{@{} l X X @{}}
\toprule
\textbf{Topología} & \textbf{Diagrama Circuital} & \textbf{Curva de Transferencia ($v_o$ vs $v_i$)} \\
\midrule
\textbf{Paralelo Positivo} & 
\begin{circuitikz}[scale=0.6, transform shape]
    \draw (0,0) to[short, o-] (1,0) to[R, l=$R$] (3,0) to[short, -o] (4,0);
    \draw (3,0) to[D] (3,-2);
    \draw (0,-2) to[short, o-o] (4,-2);
    \node[left] at (0,-1) {$v_i$}; \node[right] at (4,-1) {$v_o$};
\end{circuitikz} & 
Si $v_i > V_\gamma \implies v_o = V_\gamma$ \newline Si $v_i < V_\gamma \implies v_o = v_i$ \\

\textbf{Paralelo Negativo} & 
\begin{circuitikz}[scale=0.6, transform shape]
    \draw (0,0) to[short, o-] (1,0) to[R, l=$R$] (3,0) to[short, -o] (4,0);
    \draw (3,-2) to[D] (3,0);
    \draw (0,-2) to[short, o-o] (4,-2);
\end{circuitikz} & 
Si $v_i < -V_\gamma \implies v_o = -V_\gamma$ \newline Si $v_i > -V_\gamma \implies v_o = v_i$ \\

\textbf{Clipper Polarizado} & 
\begin{circuitikz}[scale=0.6, transform shape]
    \draw (0,0) to[short, o-] (1,0) to[R, l=$R$] (3,0) to[short, -o] (4,0);
    \draw (3,0) to[D] (3,-1.5) to[V, v=$V_{\text{ref}}$] (3,-3);
    \draw (0,-3) to[short, o-o] (4,-3);
\end{circuitikz} & 
Si $v_i > V_{\text{ref}} + V_\gamma \implies v_o = V_{\text{ref}} + V_\gamma$ \newline Si $v_i < V_{\text{ref}} + V_\gamma \implies v_o = v_i$ \\
\bottomrule
\end{tabularx}
\end{center}

% ==========================================
% SECCIÓN 3: SUJETADORES (CLAMPERS)
% ==========================================
\section{Circuitos Sujetadores de Nivel (Clampers)}
Un \textit{clamper} añade un offset de DC a una señal de CA \textbf{sin alterar su amplitud pico a pico}. Su núcleo es un capacitor que actúa como una batería flotante autocalibrable.

\vspace{0.2cm}
\noindent
\begin{minipage}{0.4\textwidth}
    \centering
    \begin{circuitikz}[american, scale=0.8, transform shape]
        \draw (0,0) to[sV, l=$v_i(t)$] (0,3) to[C, l=$C$] (3,3) to[short, -o] (5,3) node[right]{$v_o(t)$};
        \draw (3,0) to[D, l_=$D_1$] (3,3);
        \draw (5,3) to[R, l=$R_L$] (5,0);
        \draw (0,0) to[short, -o] (5,0);
    \end{circuitikz}
\end{minipage}
\begin{minipage}{0.55\textwidth}
    \textbf{Mecánica de Operación:}\\
    \textbf{1. Carga Transitoria:} En el primer semiciclo negativo, $D_1$ conduce. $C$ se carga a $V_C = V_p - V_\gamma$.\\
    \textbf{2. Régimen Permanente:} $D_1$ se apaga. $C$ se comporta como batería continua. $v_o(t) = v_i(t) + V_C$.\\
    \textbf{Condición de Diseño:} Para evitar \textit{droop} (descarga parásita), exigimos:
    \begin{equation}
        \tau = R_L C \ge 10 \cdot T
    \end{equation}
\end{minipage}

\newpage
% ==========================================
% SECCIÓN 4: PROBLEMA INTEGRADOR PFI
% ==========================================
\section{Problema de Cátedra: Acondicionamiento ADC (PFI)}
\textbf{Enunciado:} El sensor entrega $v_i(t) = 6 \sin(2\pi \cdot 1000 \cdot t)\,\text{V}$ ($V_p = \qty{6}{\volt}, f = \qty{1}{\kilo\hertz}$). El microcontrolador admite de $\qty{0}{\volt}$ a $\qty{+5.0}{\volt}$. Se diseña un circuito en dos etapas ($V_\gamma = \qty{0.7}{\volt}$):
\begin{enumerate}
    \item \textbf{Clamper:} Eleva la señal usando $R_L = \qty{47}{\kilo\ohm}$. 
    \item \textbf{Clipper:} Limita los picos a un techo seguro de $\qty{+5.1}{\volt}$ usando $R_1 = \qty{1}{\kilo\ohm}$.
\end{enumerate}

\begin{center}
\begin{circuitikz}[american, scale=0.85, transform shape]
    % Etapa 1: Clamper
    \draw (0,0) to[sV, l=$v_i(t)$] (0,3) to[C, l=$C$] (3,3);
    \draw (3,0) to[D, l_=$D_1$] (3,3);
    \draw (5,3) to[R, l=$R_L$ (\qty{47}{\kilo\ohm})] (5,0);
    \node[above] at (4,3) {$v_{o1}(t)$};
    % Etapa 2: Clipper
    \draw (3,3) -- (5,3) -- (7,3) to[R, l=$R_1$ (\qty{1}{\kilo\ohm})] (9,3);
    \draw (9,3) to[D, l=$D_2$] (9,1.5) to[V, v=$V_{\text{CLIP}}$] (9,0);
    \draw (9,3) -- (11,3) to[short, -o] (11,3) node[right]{$v_{\text{ADC}}(t)$};
    \draw (0,0) -- (11,0);
    \node[ground] at (5,0) {};
\end{circuitikz}
\end{center}

\subsection*{Resolución Analítica y Análisis Gráfico}
\textbf{Paso 1: Etapa Sujetadora}\\
Para evitar el decaimiento, $C \ge \frac{T}{0.01 \cdot R_L} = \qty{2.127}{\micro\farad} \implies \mathbf{C = \qty{2.2}{\micro\farad}}$.\\
Voltaje retenido: $V_C = 6.0 - 0.7 = \mathbf{\qty{5.3}{\volt}}$. La onda resultante se desplaza hacia arriba: $v_{o1}(t) = 6 \sin(2000\pi t) + 5.3\text{ V}$.

\textbf{Paso 2: Etapa Recortadora}\\
Para proteger el ADC, limitamos en $\qty{5.1}{\volt}$. El diodo $D_2$ entra en conducción en $V_{\text{CLIP}} + V_\gamma = 5.1 \implies V_{\text{CLIP}} = \mathbf{\qty{4.4}{\volt}}$.

\vspace{0.4cm}
\textbf{Gráfica Dinámica 1: Transitorio de Carga y Limitación en el Tiempo}\\
Observen cómo durante la primera mitad del primer ciclo el capacitor no tiene carga ($v_o = v_{in}$). Al llegar a $\qty{-0.7}{\volt}$, el capacitor absorbe energía rápidamente, fijando el offset permanente de $\qty{5.3}{\volt}$.

\begin{figure}[h!]
    \centering
    \begin{tikzpicture}
        \begin{axis}[
            width=14cm, height=6.5cm,
            axis lines=middle,
            xlabel={Tiempo ($t/T$)}, ylabel={Voltaje (\unit{\volt})},
            xmin=0, xmax=2.5, ymin=-7, ymax=13,
            grid=both, grid style={dashed, gray!30},
            legend style={at={(0.98,0.05)}, anchor=south east}
        ]
        % v_in
        \addplot[domain=0:2.5, samples=200, gray, dashed, thick] {6*sin(360*x)};
        \addlegendentry{$v_{in}(t)$ (Entrada Bipolar)}
        
        % Transitorio exacto clamper: 
        % Si x < 0.5: v_in. 
        % Si 0.5 <= x < 0.75: se clampa en -0.7 mientras carga. 
        % Si x >= 0.75: v_in + 5.3
        \addplot[domain=0:2.5, samples=400, blue, thick] {
            (x<0.5)*(6*sin(360*x)) + 
            (x>=0.5)*(x<0.75)*(max(-0.7, 6*sin(360*x))) + 
            (x>=0.75)*(max(-0.7, 6*sin(360*x) + 5.3))
        };
        \addlegendentry{$v_{o1}(t)$ (Clamper)}
        
        % v_adc: Es v_o1 limitada a 5.1V por arriba
        \addplot[domain=0:2.5, samples=400, red, ultra thick] {
            min(5.1, (x<0.5)*(6*sin(360*x)) + 
            (x>=0.5)*(x<0.75)*(max(-0.7, 6*sin(360*x))) + 
            (x>=0.75)*(max(-0.7, 6*sin(360*x) + 5.3)))
        };
        \addlegendentry{$v_{ADC}(t)$ (Acondicionada)}
        
        % Límites
        \addplot[domain=0:2.5, darkred, dashed] {5.1};
        \addplot[domain=0:2.5, black, dashed] {-0.7};
        \end{axis}
    \end{tikzpicture}
\end{figure}

\vspace{0.4cm}
\textbf{Gráfica Dinámica 2: Curva de Transferencia en Régimen Permanente}\\
Eliminando el tiempo de la ecuación, trazamos la relación directa entre la entrada y el ADC.

\begin{figure}[h!]
    \centering
    \begin{tikzpicture}
        \begin{axis}[
            width=10cm, height=6cm,
            axis lines=middle,
            xlabel={$v_{in}$ (\unit{\volt})}, ylabel={$v_{ADC}$ (\unit{\volt})},
            xmin=-7, xmax=7, ymin=-2, ymax=7,
            grid=both, grid style={dashed, gray!30},
            legend style={at={(0.05,0.95)}, anchor=north west}
        ]
        % Curva de Transferencia: min(5.1, max(-0.7, v_in + 5.3))
        \addplot[domain=-6:6, samples=200, purple, ultra thick] {min(5.1, max(-0.7, x + 5.3))};
        \addlegendentry{Curva de Transferencia $v_{ADC}$}
        
        % Puntos de inflexión
        \draw[dashed, gray] (axis cs:-6, -0.7) -- (axis cs: -6, 0);
        \draw[dashed, gray] (axis cs:-0.2, 5.1) -- (axis cs: 6, 5.1);
        \node[above] at (-0.2, 5.1) {Corte Superior};
        \node[below] at (-6, -0.7) {Corte Inferior};
        \end{axis}
    \end{tikzpicture}
\end{figure}

\newpage
% ==========================================
% SECCIÓN 5: PYTHON Y AUDITORÍA
% ==========================================
\section{Taller Computacional: Simulación Transitoria Discreta}

\begin{lstlisting}[language=Python]
# ==============================================================================
# UNIVERSIDAD CATOLICA BOLIVIANA "SAN PABLO" - SEDE TARIJA
# Sesion 6: Simulacion Transitoria de Clampers y Clippers
# Docente: M.Sc. Ing. Bernardo Quiroga Turdera
# ==============================================================================
import numpy as np
import matplotlib.pyplot as plt

frecuencia = 1000.0; T = 1.0 / frecuencia
Vp = 6.0; V_gamma = 0.7
R_L = 47e3; C = 2.2e-6
V_clip_ref = 4.4; V_max_adc = V_clip_ref + V_gamma

t = np.linspace(0, 3 * T, 4000)
dt = t[1] - t[0]
v_in = Vp * np.sin(2 * np.pi * frecuencia * t)
v_o1 = np.zeros_like(t); v_adc = np.zeros_like(t)

v_cap = 0.0; r_d_on = 10.0

for i in range(len(t)):
    v_nodo_d1 = v_in[i] - v_cap
    if v_nodo_d1 < -V_gamma: # Carga rapida a traves de D1
        i_charge = (-V_gamma - v_nodo_d1) / r_d_on
        v_cap -= (i_charge * dt) / C
        v_o1[i] = -V_gamma
    else:                    # Descarga a traves de R_L
        i_discharge = v_nodo_d1 / R_L
        v_cap -= (i_discharge * dt) / C
        v_o1[i] = v_nodo_d1
    
    # Recortador D2
    v_adc[i] = V_max_adc if v_o1[i] > V_max_adc else v_o1[i]

plt.figure(figsize=(10, 6))
plt.plot(t*1e3, v_in, 'gray', linestyle='--', label='v_in (+-6V)')
plt.plot(t*1e3, v_o1, 'b:', label='v_o1 (Clamper)')
plt.plot(t*1e3, v_adc, 'r-', linewidth=2, label='v_adc (Protegida 5.1V)')
plt.axhline(y=V_max_adc, color='darkred', linestyle='--')
plt.title('Acondicionamiento de Senal PFI'); plt.xlabel('Tiempo (ms)'); plt.ylabel('Tension (V)')
plt.grid(True); plt.legend(loc='upper right')
plt.show()
\end{lstlisting}

\begin{aviso}[Protocolo de Auditoría Dinámica (IA en Banco)]
\textbf{Consigna:} Cambien $C$ de $\qty{2.2}{\micro\farad}$ a $\qty{10}{\nano\farad}$.\\
\textbf{Pregunta:} ¿Por qué la parte inferior de la señal se deforma cayendo por debajo de $\qty{-0.7}{\volt}$?\\
\textbf{Respuesta:} Al reducir $C$, la constante $\tau = R_L C$ se vuelve mucho menor que $10T$. El capacitor se descarga masivamente perdiendo su función de ``batería constante'', ocasionando un severo \textit{droop} o decaimiento en la señal.
\end{aviso}

\end{document}